\documentclass[manuscript,anonymous]{acmart}
\setcopyright{none}
\acmYear{2027}
\acmDOI{}
\acmISBN{}
\acmConference[EICS 2027]{EICS 2027}{June 21--25, 2027}{Toulouse, France}
\settopmatter{printfolios=true}
\usepackage{booktabs}
\usepackage{array}
\hypersetup{pdfauthor={},pdftitle={Experience Engine: A Framework-Agnostic Protocol for Multi-Dimensional UI Experience Transitions}}
\title{Experience Engine: A Framework-Agnostic Protocol for Multi-Dimensional UI Experience Transitions}
\author{Anonymous Author(s)}
\affiliation{\institution{Anonymous Institution}}
\begin{document}
\begin{abstract}
Multi-dimensional web interfaces may vary simultaneously along dimensions such as culture,
theme, motion, typography, assets, resources, and component behavior. Existing technologies
address many of these dimensions individually, but application code can still be responsible
for coordinating cross-dimensional resolution, resource preparation, component adaptation,
transition readiness, and state visibility. This paper presents the Experience Transition
Protocol (ETP), a framework-agnostic runtime protocol implemented by Experience Engine.
ETP models an experience as a resolved snapshot, computes a typed experience delta, derives
the affected dependency closure, prepares heterogeneous dependent resources and component
adaptations, and performs a guarded commit. We formalize the protocol and state invariants,
implement a reference engine, and evaluate it against a strong manual Traditional+ baseline
that independently reproduces the same tested safety semantics. Across generated and
browser-based workloads, ETP preserves failure and stale-request invariants and coordinates
heterogeneous consequences across culture, theme, and motion. Ablation experiments show
that removing typed deltas, dependency closure, guarded commit, or prepare-before-commit
introduces measurable extra work or safety failures under the corresponding workloads.
The evaluation does not establish universal performance superiority, unique atomicity, or
absolute novelty; locality results are limited to the tested synthetic dependency graphs.
Exact React 18 E2E and Next.js SSR/RSC runtime validation were also performed, including
request-time cookie personalization on a production Next.js setup, while third-party replication
remains outside the present evaluation.
\end{abstract}
\keywords{interactive systems, UI engineering, adaptive interfaces, runtime systems, design systems, transition protocols}
\ccsdesc[500]{Human-centered computing~Interactive systems and tools}
\ccsdesc[300]{Software and its engineering~Software architectures}
\ccsdesc[300]{Human-centered computing~Web-based interaction}
\maketitle

\section{Introduction}
Modern UI systems increasingly adapt along multiple dimensions at runtime. A single
user-visible change can affect language and direction, design tokens, component variants,
component replacements, fonts, images, code, and motion. These concerns are often
represented by separate mechanisms, leaving application code responsible for coordinating
changes and deciding when the resulting state becomes visible.

We study a narrower systems problem: how to represent and execute a transition between
multi-dimensional UI experiences as one reusable runtime operation without coupling the
semantic ownership of culture, theme, and motion. We do not replace localization,
design-system theming, animation libraries, or framework rendering. We define a protocol
boundary at which these mechanisms can be coordinated.

The Experience Transition Protocol (ETP) makes an experience-level transition explicit,
computes a typed difference, derives a dependency closure, prepares required consequences,
and commits only an accepted result. The reference implementation, Experience Engine,
provides framework adapters above a framework-agnostic core.

The contribution is deliberately scoped. Prior work establishes substantial foundations for
adaptive UI, model-based adaptation, multi-dimensional adaptation, dependency tracking,
design-token resolution, and browser transition mechanisms. We do not claim to introduce
those ideas individually. Our claim is that ETP provides a formalized and reusable protocol
abstraction for coordinating them at the level of a multi-dimensional UI experience
transition.

\section{Problem Definition and Model}
An experience is modeled as $E=(C,T,M,K,P,R)$, where $C$, $T$, and $M$ denote culture,
theme, and motion; $K$ denotes resolved tokens and semantic properties; $P$ denotes
component decisions; and $R$ denotes dependent resources.

Semantic ownership is explicit: culture owns locale, direction, formatting, and cultural
presentation; theme owns visual tokens and theme-level component properties; motion remains
independent. Fonts and assets may be context-dependent.

The protocol distinguishes token changes, variants, and replacements. The execution sequence is
\[
\mathrm{Normalize}\rightarrow\mathrm{Resolve}\rightarrow\mathrm{Diff}\rightarrow
\mathrm{Seeds}\rightarrow\mathrm{Closure}\rightarrow\mathrm{Prepare}\rightarrow
\mathrm{Transition}\rightarrow\mathrm{Commit}.
\]


\begin{figure}[t]
\centering
\fbox{\begin{minipage}{0.92\columnwidth}
\centering\textbf{Experience Transition Protocol}\\[3pt]
\footnotesize
Normalize $\rightarrow$ Resolve $\rightarrow$ Diff $\rightarrow$ Seeds $\rightarrow$
Closure $\rightarrow$ Prepare $\rightarrow$ Transition $\rightarrow$ Guarded Commit\\[5pt]
\begin{tabular}{c|c|c}
\textbf{Input} & \textbf{Coordination} & \textbf{Visibility}\\
Target request & Delta + closure + preparation & Guarded commit
\end{tabular}
\end{minipage}}
\Description{The Experience Transition Protocol moves from normalization and resolution through typed differences, dependency closure, preparation, transition, and guarded commit.}
\caption{ETP transition pipeline and the guarded visibility boundary.}
\label{fig:etp-pipeline}
\end{figure}

\section{Protocol and Invariants}
ETP defines a transition boundary around an immutable current snapshot and a proposed target
snapshot. A typed delta identifies directly affected dimensions. Seed computation finds direct
changes, dependency closure expands the affected work set, preparation establishes readiness
without replacing the committed snapshot, and a guarded commit accepts only the latest valid
prepared result.

The evaluated invariants cover deterministic resolution, immutable snapshots, failure
preservation, stale dominance, atomic visibility, dependency closure, semantic ownership,
and localized preparation under the tested dependency model. These invariants are not claimed
to be unique to ETP; the Traditional+ baseline reproduces the same tested safety semantics.

\section{Reference Implementation}
Experience Engine separates a framework-agnostic core from React and Next.js adapters.
The React-facing design uses an external-store integration and exposes provider and hook APIs.
The inspected core source boundary contains no direct browser-global dependency.

Registries remain declarative and resolution results immutable. Preload acquires resources,
whereas prepare establishes transition readiness. The public transition operation accepts a
complete target request and coordinates normalization, resolution, delta computation,
dependency closure, preparation, transition, and guarded commit.

\section{AI-Assisted Research Disclosure}
Generative AI tools were used as research assistants during development of the artifact,
including parts of research planning, architectural exploration, code drafting, test generation,
experiment orchestration, data-analysis assistance, validation checks, and preparation of
reproducibility materials. Human authors remain responsible for verifying the implementation,
experimental results, references, claims, and final submission. The human author should
verify that this disclosure exactly describes the actual workflow.

\section{Independent Reference-Model Reproduction}
A separate semantic reference implementation was written without importing Experience Engine.
It independently implemented Normalize, Resolve, Diff, Seeds, Closure, Prepare, and Commit.
The model passed 1,000 identity checks and 5,000 randomized transitions, including
deterministic failure injection, with zero observed failure-preservation violations. This is
an author-created independent implementation rather than third-party replication.


\begin{table}[t]
\centering
\caption{Final empirical evidence summary. Results are scoped to the reported workloads.}
\label{tab:evidence}
\small
\begin{tabular}{p{0.24\columnwidth}p{0.18\columnwidth}p{0.45\columnwidth}}
\toprule
Evidence & Size & Outcome\\
\midrule
Traditional+ & 10,000 & 0 partial commits; 0 stale overwrites\\
Locality DAG & 10,000 nodes / 6 required & 99.94\% reduction vs full invalidation\\
Heterogeneous & 5,000 random & 0 semantic mismatches\\
Browser stress & 301 cases & 0 invariant violations\\
React 18 E2E & 3 transitions & 3/3 state matches; 3/3 React re-renders\\
Next dynamic SSR & 2 cookie cases & 2/2 exact server identities; HTTP 200\\
Ablation & 10,000 transitions & Stage-specific failures/work effects\\
Reference model & 5,000 random & 0 failure-preservation violations\\
\bottomrule
\end{tabular}
\end{table}

\section{Evaluation}
The evaluation combines formal/property-style testing, a separate reference-model
reproduction, synthetic dependency workloads, heterogeneous adaptation workloads, browser
stress tests, and a manual Traditional+ baseline.

\subsection{Verification}
Generated tests covered deterministic identity, immutable snapshots, direction handling,
failure preservation, stale dominance, and dependency closure. No invariant violations
were observed in the stored generated workloads.

\subsection{Traditional+ Baseline}
The strong manual baseline contains a central transaction controller, explicit dependency
graph, semantic delta, prepare-before-commit, generation guards, and manual
resource/component/motion coordination. Across 10,000 matched transitions with injected
resource failures and stale-request scenarios, ETP and Traditional+ were semantically
equivalent in the tested safety outcomes: no partial commits and no stale overwrites.

\subsection{Locality}
Structured dependency DAG experiments evaluated 1,000, 5,000, and 10,000-node graphs.
In the strongest 10,000-node case with six required nodes, ETP prepared the six-node closure
rather than invalidating all 10,000 nodes, a 99.94\% reduction relative to full invalidation.
These results characterize the tested synthetic dependency model.

\subsection{Heterogeneous Consequences}
Across 12 combinations of two cultures, two themes, and three motion modes, all 144 ordered
transitions were exercised, followed by 5,000 random transitions. Consequences included
tokens, variants, replacements, fonts, translations, assets, and motion. Semantic comparison
against the strong manual baseline found 0 mismatches in the stored random workload.

\subsection{Browser Stress}
Real headless Chromium execution exercised 50 resource-failure cases, 50 stale-race cases,
101 rapid requests, and 100 mixed cases. The stored results showed zero invariant violations.
In the rapid-request scenario, one request committed and 100 earlier requests were classified
stale, with the final state matching the final request.

\subsection{Framework Runtime Validation}
The implementation was validated through an exact React 18.3.1 browser path and a Next.js 16.2.0 App Router path. In React, the test mounted `ExperienceProvider` with `createRoot`, used `useExperience` and `useSyncExternalStore`, executed three experience transitions through the engine, and verified that requested and final experience states matched in all three cases; an actual React re-render was observed for each transition. The observed switch latencies were 23.9, 34.0, and 27.4 ms, with CLS 0 and zero dropped frames in the stored samples. Long-task observation reported two entries per sample, so these measurements are treated as characterization rather than a performance claim.

For Next.js, the production build completed successfully using App Router. The initial server response was verified at HTTP 200 and contained the server-component markers for `ar-EG`, `luxury`, `smooth`, and RTL direction. A client-side benchmark then executed three transitions through the real React adapter and verified DOM convergence for each request. Finally, a dynamic route at `/personalized` was exercised with request-time `experience` cookies. The same route returned distinct server-resolved identities: `ar-EG::luxury::smooth` for `experience=ar-EG.luxury.smooth` and `en-US::light::instant` for `experience=en-US.light.instant`, both with HTTP 200 and matching direction. This demonstrates request-time personalization in the tested Next.js production setup; it does not generalize to every deployment or cache topology.

\subsection{Performance}
Core microbenchmarks do not support a universal performance-superiority claim. In several
measured scenarios, the engine introduced orchestration overhead. Performance is therefore
treated as systems characterization.

\subsection{Independent Reference-Model Reproduction}
A separate semantic reference implementation was written without importing Experience Engine.
It independently implemented Normalize, Resolve, Diff, Seeds, Closure, Prepare, and Commit.
The model passed 1,000 identity checks and 5,000 randomized transitions, including deterministic
failure injection, with zero observed failure-preservation violations. This is an
author-created independent implementation rather than third-party replication.

\subsection{Final Evidence Audit}
The final audit examined 11 project claims. The primary retained claims are the formalized
ETP protocol abstraction and its coordination of typed delta, dependency closure,
heterogeneous preparation, and guarded commit. Claims of unique safety semantics, universal
performance superiority, and being the first adaptive or multi-dimensional UI framework were
rejected. Locality and browser results remain explicitly scoped. React 18 E2E and Next.js
SSR/RSC runtime validation are supported by direct runtime evidence within the tested setups;
third-party replication remains pending. Arithmetic consistency checks passed.

\section{Ablation}
The ablation study used 500 generated dependency graphs and 10,000 transitions with 12\%
failure injection and 12\% stale-request injection. Removing typed delta increased the
prepared-work ratio to 1.358$\times$. Removing dependency closure caused 7,661 downstream
omissions. Removing the guard caused 1,050 stale overwrites. Removing prepare-before-commit
caused 1,200 partial states.


\begin{table}[t]
\centering
\caption{Final claim boundary.}
\label{tab:claim-boundary}
\small
\begin{tabular}{p{0.28\columnwidth}p{0.18\columnwidth}p{0.40\columnwidth}}
\toprule
Claim & Status & Treatment\\
\midrule
ETP protocol abstraction & Supported & Primary contribution\\
Typed delta + closure + preparation + commit & Supported & Core protocol\\
Synthetic locality & Scoped & Workload-specific\\
Browser stress & Scoped & Tested Chromium scenarios\\
Unique safety semantics & Rejected & Traditional+ reproduces tested semantics\\
Universal speedup & Rejected & Performance is characterization\\
First adaptive UI & Rejected & Prior-art audit does not support it\\
React 18 exact E2E & Supported & Direct runtime evidence\\
Next SSR/RSC & Supported & Direct runtime evidence within tested setup\\
Third-party replication & Pending & External replication still required\\
\bottomrule
\end{tabular}
\end{table}

\section{Threats to Validity and Scope}
The strongest tested safety semantics are reproducible by Traditional+. Locality is based on
structured synthetic graphs. Browser evidence is scoped to controlled headless Chromium
workloads. Exact React 18 E2E and Next.js SSR/RSC runtime validation were completed in the local test environments. The Next.js dynamic per-request test used request cookies on a production server and therefore does not establish behavior for every deployment, CDN, or cache topology.
Artifact-level reproducibility and author-created independent reference modeling are not
third-party replication. The work evaluates system properties rather than user outcomes.

The prior-art audit did not establish absolute novelty over adaptive UI, model-based
adaptation, multi-dimensional adaptation, dependency tracking, or transactional UI
techniques. The contribution claim is intentionally limited to the ETP abstraction and its
explicit protocol framing.

\section{Related Work}
Adaptive and model-based UI research has long studied runtime transformation, context-aware
adaptation, multi-dimensional adaptation, and mappings between abstract and concrete UI
realizations \cite{nilsson2006,calvary2003}. Culturally adaptive UI work studies the
relationship between culture and interface presentation \cite{reinecke2011}. Design-token resolution and browser view transitions provide mature
building blocks for visual and transition concerns \cite{dtcgresolver,w3cvt}.

ETP is positioned as a coordination layer across these mechanisms. It does not replace
localization, design systems, animation primitives, or browser transition APIs. Its scope
is the experience-level protocol that makes the transition itself a first-class runtime
operation, including typed change detection, dependency closure, heterogeneous preparation,
and guarded acceptance.

\section{Final Claim Boundary}
The final evidence audit retains two primary claims: ETP is a formalized and reusable
protocol abstraction for experience-level state transitions, and it coordinates typed
delta, dependency closure, heterogeneous preparation, and guarded commit. Claims of
unique safety semantics, universal performance superiority, and priority over the
adaptive-UI literature are excluded. React 18 E2E and Next.js SSR/RSC runtime validation
are supported by direct runtime evidence within the tested setups; third-party replication
remains pending.

\section{Conclusion}
ETP is a formalized and reusable protocol abstraction for experience-level state transitions
in multi-dimensional UI systems. The evidence supports the abstraction within the tested
scope: formal properties pass, heterogeneous adaptation matches the strong manual baseline,
browser stress tests preserve tested safety invariants, ablation exposes the role of
protocol stages, and an independent reference model reproduces the core semantic behavior.

The evidence also constrains the claim: manual code can reproduce the tested safety semantics,
the engine is not universally faster, synthetic locality results are workload-specific, and
third-party replication remains open; the Next.js dynamic SSR evidence is scoped to the tested production configuration.

\begingroup\scriptsize
\begin{thebibliography}{9}
\setlength{\itemsep}{0pt}\setlength{\parsep}{0pt}\setlength{\parskip}{0pt}
\bibitem{nilsson2006}
Erik G. Nilsson, Jacqueline Floch, Svein O. Hallsteinsen, and Erlend Stav.
2006. Model-based user interface adaptation. \emph{Computers \& Graphics} 30, 5 (2006), 692--701.
\bibitem{reinecke2011}
Katharina Reinecke and Abraham Bernstein. 2011. Improving Performance, Perceived Usability,
and Aesthetics with Culturally Adaptive User Interfaces. \emph{ACM TOCHI} 18, 2, Article 8.
\bibitem{calvary2003}
Gaelle Calvary, Joelle Coutaz, Olivier Thevenin, and others. 2003.
A unifying reference framework for multi-target user interfaces.
\bibitem{dtcgresolver}
Design Tokens Community Group. 2025. Design Tokens Resolver Module 2025.10.
\bibitem{w3cvt}
W3C CSS Working Group. CSS View Transitions Module Level 1.
\end{thebibliography}
\endgroup
\end{document}
